\section{Conclusiones}

El software ICOCIV constituye una herramienta técnica de apoyo para el análisis y proyección de índices de costos de obras civiles publicados por el DANE. Su valor no reside únicamente en automatizar cálculos, sino en organizar un flujo metodológico completo: lectura robusta de datos oficiales, validación temporal, selección jerárquica, modelación interpretable, diagnóstico estadístico, backtesting y generación de informes técnicos.

Desde el punto de vista estadístico, el documento muestra que las series ICOCIV deben tratarse como series temporales económicas. No son observaciones IID normales. Presentan dependencia temporal, tendencia, posible autocorrelación y sensibilidad a condiciones económicas. Por esta razón, el software incorpora validación cronológica y backtesting temporal en lugar de evaluaciones aleatorias que romperían la estructura pasado-futuro.

La regresión lineal, polinómica, logarítmica y robusta ofrece un conjunto de modelos suficientemente flexible para representar trayectorias comunes de índices, sin sacrificar interpretabilidad. La selección automática no debe basarse solo en \(R^2\), sino en un balance entre ajuste, parsimonia, residuos, AICc, estabilidad predictiva e incertidumbre. Esta lógica es coherente con la literatura de forecasting y econometría aplicada.

El análisis residual fortalece la confiabilidad del sistema porque permite identificar patrones no explicados, autocorrelación, desviaciones de normalidad y posibles problemas de especificación. A su vez, los intervalos de confianza o predicción recuerdan que toda proyección contiene incertidumbre. Una estimación futura debe interpretarse como soporte cuantitativo condicionado, no como certeza contractual.

El backtesting temporal es uno de los elementos metodológicos más importantes del software. Al simular predicciones históricas usando únicamente información pasada, permite evaluar capacidad predictiva real. Métricas como MAE, RMSE y MAPE ofrecen una lectura práctica del error y ayudan a justificar la elección del modelo ante usuarios técnicos y jurados académicos.

En términos de ingeniería civil, la interpretación estadística debe complementarse con contexto económico. Materiales, transporte, combustibles, mano de obra, inflación y condiciones internacionales pueden influir en los índices. El software aporta evidencia numérica, pero el juicio profesional sigue siendo necesario para valorar riesgos, choques externos y pertinencia de uso.

La decisión de evitar modelos de caja negra en esta etapa es metodológicamente defendible. Para un sistema institucional, local y orientado a documentación técnica, la transparencia es prioritaria. Los modelos interpretables permiten explicar supuestos, limitaciones y resultados con claridad, lo cual fortalece la reproducibilidad y la aceptación académica del sistema.

Finalmente, el software representa una evolución desde una calculadora de proyecciones hacia un sistema de análisis estadístico y validación predictiva de índices ICOCIV. Su aporte académico consiste en integrar datos oficiales, fundamentos econométricos y documentación reproducible para apoyar el análisis de costos en obras civiles. Su aporte técnico consiste en ofrecer una base local, liviana y trazable para producir informes que puedan respaldar decisiones profesionales e institucionales.
